%% ch11.tex — Chapter 11: The Lakehouse, MLOps, and Emerging Architectures
%% Part of "Data Engineering" textbook

\chapter{The Lakehouse, MLOps, and Emerging Architectures}
\label{ch:modern-stack}

\epigraph{\textit{``The exciting thing about the Lakehouse is that it is
not a compromise --- it is actually better than either a data warehouse
or a data lake at their own jobs.''}}{Matei Zaharia,
co-creator of Apache Spark and Delta Lake}

\vspace{4pt}

\begin{chapterlearning}
After completing this chapter, you will be able to:
\begin{enumerate}
  \item Describe the four eras of big data technology (2003--present),
        explain the engineering problem that drove each transition,
        and identify which era's infrastructure remains relevant in
        modern production deployments.
  \item Explain the Lakehouse architecture: how an open table format
        (Delta Lake, Apache Iceberg, Apache Hudi) stored on cheap object
        storage achieves ACID transactions via optimistic concurrency
        control, time travel via a transaction log, and schema enforcement
        at the storage layer.
  \item Compare Delta Lake's transaction log mechanism against the
        traditional data warehouse approach and the raw data lake approach
        on dimensions of ACID compliance, open format, BI support, ML support,
        time travel, and storage cost.
  \item Explain Apache Iceberg's hidden partitioning and partition
        evolution capabilities and identify the scenario where Iceberg
        is preferred over Delta Lake.
  \item Describe the four principles of the Data Mesh paradigm (domain
        ownership, data as a product, self-serve platform, federated
        computational governance) and explain which engineering bottleneck
        each principle addresses.
  \item Identify the seven components of an MLOps pipeline (ingestion,
        feature engineering, feature store, model training, experiment
        tracking, model registry, serving and monitoring), describe the
        tool used for each, and explain what ``hidden technical debt'' each
        component eliminates.
  \item Describe the three-tier edge--cloud architecture for IoT analytics
        (device, edge gateway, cloud) and explain what processing happens
        at each tier and why.
  \item Define $(\varepsilon, \delta)$-differential privacy formally,
        explain the Laplace noise mechanism, and describe the accuracy--
        privacy trade-off controlled by the privacy budget $\varepsilon$.
  \item Describe the Dataflow Model's four questions (What, Where, When,
        How), explain how batch processing is a special case of streaming
        under this model, and explain how Apache Beam eliminates the Lambda
        architecture's dual code path.
  \item Apply Delta Lake's core operations (create, append, update, time
        travel, DESCRIBE HISTORY) using PySpark and the \code{delta-spark}
        library.
\end{enumerate}
\end{chapterlearning}

\bigskip

Every chapter in this book has described a layer of the data engineering
stack as it existed at a specific moment in the field's development.
Chapter~\ref{ch:hdfs} described HDFS and the assumption that storage
and compute were co-located on-premise. Chapter~\ref{ch:mapreduce}
described MapReduce's batch processing model. Chapter~\ref{ch:spark}
described how Spark replaced MapReduce with in-memory processing.
Chapter~\ref{ch:pipelines} described Kafka, streaming pipelines, and the
Lambda architecture's attempt to unify batch and stream.

This final chapter completes the arc. It describes the fourth era of
big data technology---the \term{Lakehouse}---which resolves the central
tension between data lakes (cheap, flexible, but unreliable) and data
warehouses (reliable, fast, but expensive and closed). It also covers
the organisational architecture (\term{Data Mesh}) that resolves the
bottleneck created when a single platform team owns all data pipelines;
the operational discipline (\term{MLOps}) that bridges the gap between
a trained model and a deployed, monitored production system; the
computational pattern (\term{edge analytics}) that extends distributed
processing to IoT devices; and the theoretical framework
(the \term{Dataflow Model}) that eliminates the Lambda architecture's
fundamental complexity by treating batch as a special case of streaming.

%%====================================================================
\section{The Technology Evolution Arc}
\label{sec:evolution}
%%====================================================================

Big data technology has evolved through four identifiable eras, each
driven by a specific engineering bottleneck that the previous era's tools
could not resolve.

\begin{table}[htbp]
\centering
\renewcommand{\arraystretch}{1.5}
\begin{tabular}{p{1.2cm}p{1.8cm}p{3.0cm}p{3.0cm}p{2.6cm}}
  \toprule
  \textbf{Era} & \textbf{Period}
    & \textbf{Defining Tools} & \textbf{Key Innovation}
    & \textbf{Bottleneck Resolved} \\
  \midrule
  1 & 2003--2011
    & GFS, MapReduce, Hadoop, Hive
    & Commodity hardware cluster replaces supercomputer
    & Storage cost; single-machine limit \\
  2 & 2011--2017
    & Spark, Kafka, YARN, Parquet
    & In-memory processing; persistent event log
    & MapReduce latency; batch-only processing \\
  3 & 2017--2022
    & AWS EMR, S3, Dataproc, Snowflake
    & Compute and storage separation; managed services
    & On-premise hardware ops; cluster idle cost \\
  4 & 2022--now
    & Delta Lake, Iceberg, MLflow, Apache Beam
    & ACID on object stores; unified BI+ML; batch=stream
    & Data lake unreliability; Lambda complexity; ML ops debt \\
  \bottomrule
\end{tabular}
\caption{Four eras of big data technology. Tools from earlier eras are
         still running in production at large enterprises --- technologies
         stratify into legacy tiers rather than disappearing.}
\label{tab:eras}
\end{table}

\begin{keyinsight}{Technologies Stratify, They Do Not Die}
HDFS and MapReduce from Era 1 are still running in production at
legacy installations across banking, telecommunications, and government.
A big data architect in 2025 must understand all four eras because they
will encounter systems from each era in real deployments. Era 1 systems
have often been running for 10+ years and cannot be migrated without
significant business risk. The progression in Table~\ref{tab:eras} is
not a replacement sequence---it is an accretion sequence. Each era added
a new layer without fully eliminating the previous one.
\end{keyinsight}

%%====================================================================
\section{Lakehouse Architecture}
\label{sec:lakehouse}
%%====================================================================

The Lakehouse resolves the central architectural tension of the 2010s:
organisations needed both the flexibility and cost-efficiency of a data
lake \emph{and} the reliability and query performance of a data warehouse,
but running both simultaneously doubled infrastructure cost and created
data consistency problems.

%--------------------------------------------------------------------
\subsection{The Two-System Problem}
\label{subsec:two-system}
%--------------------------------------------------------------------

By 2018, the standard enterprise architecture for big data analytics
comprised two separate systems:

\textbf{The data lake} (HDFS or Amazon S3 with Spark): raw data in
Parquet or ORC format, accessible to ML frameworks (Spark MLlib, PyTorch)
for model training. Cheap storage ($\sim$\$23/TB-month on S3). No ACID
guarantees: concurrent writes could corrupt data; there was no schema
enforcement to prevent malformed records from entering the lake; and
deleting a specific record (required by GDPR's right to erasure) required
reading and rewriting every Parquet file containing that record.

\textbf{The data warehouse} (Amazon Redshift, Google BigQuery, Snowflake):
ACID-compliant SQL tables, fast query performance, schema enforcement.
Expensive proprietary storage ($\sim$\$200--250/TB-month). Closed format:
ML frameworks could not read warehouse data directly, so data had to be
extracted as CSV or Parquet and loaded into the lake---a nightly ETL
copy that created a 12-hour freshness gap and a second copy of all data.

The result was three compounding problems:
\begin{enumerate}
  \item \textbf{Cost}: two systems, two storage bills, two compute bills.
  \item \textbf{Consistency}: analysts queried the warehouse; data scientists
        queried the lake. They worked with different copies of the data
        and frequently produced inconsistent results.
  \item \textbf{Lambda complexity}: streaming workloads required maintaining
        two code paths (one for the lake's batch processing, one for the
        real-time layer) that had to be kept in sync.
\end{enumerate}

%--------------------------------------------------------------------
\subsection{Delta Lake: ACID on Object Stores}
\label{subsec:delta-lake}
%--------------------------------------------------------------------

\autocite{DeltaLake2020} introduced Delta Lake as an open-source storage
layer that adds ACID transactions, scalable metadata handling, and audit
history to data lakes on object stores. Its core mechanism is the
\term{Delta Log}: a series of JSON files stored inside the table's
directory in S3 that form a serialisable transaction log.

\begin{defbox}{Delta Lake Transaction Log (Delta Log)}
A \textbf{Delta table} on S3 consists of two components stored in the
same S3 prefix:
\begin{enumerate}
  \item \textbf{Data files}: immutable Parquet files containing the actual
        row data. They are never modified in place; instead, new files are
        written and old files are logically deleted.
  \item \textbf{Delta Log} (\texttt{\_delta\_log/}): a sequence of JSON
        files (one per committed transaction: \texttt{00000.json},
        \texttt{00001.json}, \ldots) each recording the set of files
        \emph{added} and \emph{removed} by that transaction, plus schema
        changes and metadata. Every 10 commits, a \emph{checkpoint} Parquet
        file condenses the full current state of the table for fast reader
        initialisation.
\end{enumerate}
\textbf{ACID properties:}
\begin{itemize}
  \item \textbf{Atomicity}: S3's PUT operation is atomic. A transaction
        either appends its JSON entry completely or not at all. A partial
        write has no effect because readers only follow committed log entries.
  \item \textbf{Isolation (Optimistic Concurrency Control)}: writers read
        the latest committed version, make their changes, and attempt to
        write the next log entry. If two writers simultaneously attempt
        to write the same log entry number, S3's conditional PUT ensures
        only one succeeds; the other detects the conflict and either retries
        or fails with a \texttt{ConcurrentModificationException}.
  \item \textbf{Consistency}: the Delta Log is the single source of truth.
        Readers always reconstruct the latest committed table state from
        the log before scanning data files; they can never see a partial
        write.
  \item \textbf{Durability}: all committed log entries persist in S3
        indefinitely. No in-memory coordinator holds state; the log is
        self-contained.
\end{itemize}
\end{defbox}

\textbf{Time travel.} Because every version of the table is preserved
in the Delta Log, Delta Lake supports querying any past snapshot:

\begin{lstlisting}[style=hiveql, caption={Delta Lake time travel queries}]
-- Read table as of a specific version
SELECT * FROM transactions VERSION AS OF 42;

-- Read table as of a specific timestamp
SELECT *
FROM   transactions
TIMESTAMP AS OF '2025-01-01 00:00:00';

-- Inspect the full transaction history
DESCRIBE HISTORY transactions;

-- Roll back to a previous version (overwrites current version)
RESTORE TABLE transactions TO VERSION AS OF 42;
\end{lstlisting}

Time travel enables: (a)~compliance audit trails (``what did the data look
like on December 31?''); (b)~reproducible ML experiments (lock training
data to a specific version, ensuring identical results on rerun);
(c)~data recovery from accidental deletes or overwrites.

\textbf{Schema enforcement and evolution.} Delta Lake rejects any write
that does not conform to the table's declared schema---including new
columns, missing columns, or type mismatches---unless schema evolution
is explicitly enabled with the \code{mergeSchema} option. Schema enforcement
eliminates the ``data swamp'' failure mode where successive writes from
different teams gradually make the table's schema inconsistent.

\textbf{Z-Ordering.} Delta Lake supports \term{Z-ordering}: a data
locality optimisation that co-locates rows with similar values in the
same Parquet files, reducing the number of files scanned for multi-column
filter queries. A traditional Hive table partitioned by \texttt{date}
still requires scanning all files for a query that filters by
\texttt{(user\_id, date)}. Z-ordering on \texttt{(user\_id, date)}
co-locates the data so that up to 95\% of files can be skipped.

\begin{lstlisting}[style=pyspark, caption={Delta Lake core operations in PySpark}]
from pyspark.sql import SparkSession
from delta import configure_spark_with_delta_pip
from pyspark.sql.functions import col

# --- Session setup ---
builder = (SparkSession.builder
    .appName("DeltaLakeDemo")
    .config("spark.sql.extensions",
            "io.delta.sql.DeltaSparkSessionExtension")
    .config("spark.sql.catalog.spark_catalog",
            "org.apache.spark.sql.delta.catalog.DeltaCatalog"))

spark = configure_spark_with_delta_pip(builder).getOrCreate()

TABLE_PATH = "hdfs:///lakehouse/payments/transactions"

# --- 1. Create initial Delta table ---
from pyspark.sql.types import StructType, StructField, \
    StringType, DoubleType, TimestampType

schema = StructType([
    StructField("txn_id",     StringType(),    False),
    StructField("sender",     StringType(),    True),
    StructField("receiver",   StringType(),    True),
    StructField("amount_usd", DoubleType(),    True),
    StructField("ts",         TimestampType(), True),
    StructField("region",     StringType(),    True),
])

df_initial = spark.createDataFrame([
    ("T001", "user_a1b2c3", "user_d4e5f6", 500.0,  None, "Northeast"),
    ("T002", "user_g7h8i9", "user_j0k1l2", 1200.0, None, "West"),
], schema=schema)

(df_initial.write
    .format("delta")
    .mode("overwrite")
    .partitionBy("region")
    .save(TABLE_PATH))

# --- 2. Append new batch (transaction 2  --  a new commit) ---
df_new = spark.createDataFrame([
    ("T003", "user_m3n4o5", "user_a1b2c3", 250.0, None, "Northeast"),
    ("T004", "user_p6q7r8", "user_s9t0u1", 750.0, None, "South"),
], schema=schema)

(df_new.write
    .format("delta")
    .mode("append")
    .save(TABLE_PATH))

# --- 3. UPDATE  --  raise a specific amount ---
spark.sql(f"""
    UPDATE delta.`{TABLE_PATH}`
    SET    amount_usd = amount_usd * 1.10
    WHERE  txn_id = 'T001'
""")

# --- 4. DELETE  --  GDPR right to erasure ---
spark.sql(f"""
    DELETE FROM delta.`{TABLE_PATH}`
    WHERE  sender = 'user_g7h8i9'
""")

# --- 5. Time travel: read state before DELETE ---
v2 = (spark.read
    .format("delta")
    .option("versionAsOf", 2)   # version 0=create, 1=append, 2=update
    .load(TABLE_PATH))
print("State at version 2 (before DELETE):")
v2.show()

# --- 6. Inspect transaction history ---
spark.sql(f"DESCRIBE HISTORY delta.`{TABLE_PATH}`").show(truncate=False)

# --- 7. Z-ordering for faster multi-column queries ---
spark.sql(f"""
    OPTIMIZE delta.`{TABLE_PATH}`
    ZORDER BY (sender, region)
""")

# --- 8. Schema enforcement: this will FAIL with AnalysisException
try:
    bad = spark.createDataFrame([(999,)], ["txn_id"])   # wrong type
    bad.write.format("delta").mode("append").save(TABLE_PATH)
except Exception as e:
    print(f"Schema enforcement blocked write: {type(e).__name__}")
\end{lstlisting}

%--------------------------------------------------------------------
\subsection{Apache Iceberg and Apache Hudi}
\label{subsec:iceberg-hudi}
%--------------------------------------------------------------------

Delta Lake is not the only open table format. Two alternatives address
different production requirements.

\textbf{Apache Iceberg} (Netflix, open-sourced 2020) uses a \emph{manifest
list} approach to the transaction log---a tree of JSON and Avro files
that records which data files belong to each snapshot. Iceberg's key
differentiating features are:
\begin{itemize}
  \item \textbf{Hidden partitioning}: partition columns are computed
        automatically from data columns and stored in table metadata,
        not exposed in SQL queries. A table partitioned by
        \texttt{months(event\_time)} requires no changes to queries
        when the partition scheme changes---contrast with Hive's explicit
        partition column requirement.
  \item \textbf{Partition evolution}: the partitioning scheme can change
        without rewriting existing data. Old partitions are read under the
        old scheme; new data uses the new scheme. This is impossible in Hive.
  \item \textbf{Multi-engine support}: Iceberg tables can be read by
        Spark, Trino, Flink, Dremio, and Presto simultaneously without
        format conversion. Delta Lake (as of 2024) requires a compatibility
        layer (UniForm) for non-Spark engines.
\end{itemize}

\textbf{Apache Hudi} (Uber, open-sourced 2019) is optimised for
\emph{record-level upserts}: updating or deleting individual rows in a
large Parquet dataset without rewriting entire files. Hudi is the natural
choice when the primary workload is CDC (Change Data Capture)---streaming
database change events from MySQL or PostgreSQL into the lake and applying
them as upserts in near-real time.

\begin{table}[htbp]
\centering
\renewcommand{\arraystretch}{1.4}
\begin{tabular}{lp{3.2cm}p{3.2cm}p{3.2cm}}
  \toprule
  \textbf{Feature} & \textbf{Delta Lake}
    & \textbf{Apache Iceberg} & \textbf{Apache Hudi} \\
  \midrule
  Origin & Databricks & Netflix & Uber \\
  Transaction log & JSON files & Manifest list (Avro) & Timeline (Avro) \\
  Best for & Spark-native workloads & Multi-engine; partition evolution
    & Record-level CDC upserts \\
  Hidden partitioning & No & Yes & Partial \\
  Time travel & Yes & Yes & Yes \\
  Partition evolution & Limited & Full & Partial \\
  Z-Ordering / Sort & Yes & Sort order (Puffin) & Clustering \\
  Primary ecosystem & Databricks / Spark & Trino / Spark / Flink
    & Spark / Flink \\
  \bottomrule
\end{tabular}
\caption{Comparison of the three major open table formats. Delta Lake
         is the default choice for Spark-centric shops; Iceberg is
         preferred when multiple query engines must read the same tables;
         Hudi is preferred for high-throughput CDC upsert workloads.}
\label{tab:table-formats}
\end{table}

\begin{defbox}{Lakehouse Architecture (Zaharia et al., CIDR 2021)}
A \textbf{Lakehouse} is a data management system with the following
properties:
\begin{enumerate}
  \item \textbf{Cheap, open-format storage}: data is stored as Parquet
        or ORC files on commodity object storage (S3, Azure Data Lake
        Storage, Google Cloud Storage) at $\sim$\$23/TB-month.
  \item \textbf{Metadata and governance layer}: an open table format
        (Delta Lake, Iceberg, or Hudi) adds ACID transactions, schema
        enforcement, time travel, and audit logging on top of the
        object storage---without moving the data to a proprietary format.
  \item \textbf{Unified query engine}: a single query engine (Spark SQL,
        Trino, or DuckDB) serves both BI dashboards (interactive SQL)
        and ML training jobs (DataFrame reads) from the \emph{same ACID-
        consistent copy} of the data. There is no ETL copy between a
        lake and a warehouse.
  \item \textbf{Declarative access control}: a unified catalogue and
        governance layer (Databricks Unity Catalog, Apache Atlas,
        AWS Glue Data Catalog) applies column-level masking, row-level
        filtering, and audit logging automatically for every query.
\end{enumerate}
\end{defbox}

%%====================================================================
\section{Data Mesh: Decentralised Domain Ownership}
\label{sec:data-mesh}
%%====================================================================

The Lakehouse solves the technical problem of unreliable, expensive,
fragmented data storage. The Data Mesh \autocite{Dehghani2022} addresses
a different class of problem: the \emph{organisational} bottleneck created
when a central data platform team is responsible for serving every business
domain's data needs.

%--------------------------------------------------------------------
\subsection{The Central Platform Bottleneck}
\label{subsec:mesh-problem}
%--------------------------------------------------------------------

In the traditional model, a central data engineering team:
\begin{enumerate}
  \item Ingests raw data from every business domain (sales, finance,
        operations, marketing) into a centralised data lake.
  \item Builds transformation pipelines to produce clean, curated datasets.
  \item Publishes datasets for consumption by analysts and data scientists.
\end{enumerate}

At 200+ teams in a large enterprise, this creates a predictable bottleneck:
\begin{itemize}
  \item \textbf{Ownership mismatch}: the pipeline team does not understand
        the domain semantics (what does a ``cancelled order'' mean in the
        sales domain?); the domain team that understands the data cannot
        fix the pipeline.
  \item \textbf{Queue bottleneck}: 200 teams compete for 20 data engineers.
        New dataset requests take weeks or months to fulfil.
  \item \textbf{Quality degradation}: no single team is accountable for
        downstream data correctness. When a dashboard shows incorrect
        numbers, three teams blame each other.
\end{itemize}

%--------------------------------------------------------------------
\subsection{The Four Data Mesh Principles}
\label{subsec:mesh-principles}
%--------------------------------------------------------------------

Dehghani proposed Data Mesh as a \emph{sociotechnical} paradigm---meaning
it changes both the technology architecture and the organisational structure
simultaneously. It has four principles.

\begin{defbox}{Data Mesh: Four Principles (Dehghani, 2019/2022)}
\begin{enumerate}
  \item \textbf{Domain ownership}: each business domain (sales, finance,
        logistics, marketing) owns, builds, operates, and is SLA-accountable
        for its own data products. The sales team is responsible for the
        quality of \texttt{sales\_events}; the logistics team is responsible
        for \texttt{shipment\_records}. There is no central data team that
        intermediates.
  \item \textbf{Data as a product}: each domain's dataset is treated as
        a product with explicit quality contracts. A \emph{data product}
        is \textbf{discoverable} (findable in a catalogue), \textbf{addressable}
        (accessible via a stable, versioned API or URI), \textbf{trustworthy}
        (with documented SLAs on freshness, accuracy, and completeness),
        and \textbf{self-describing} (with schema, lineage, and ownership
        metadata).
  \item \textbf{Self-serve data platform}: domain teams are not data
        engineers. The central platform team's role changes from
        \emph{owning pipelines} to \emph{providing infrastructure
        abstractions} so that domain teams can build data products without
        Spark or Kafka expertise. The Lakehouse is the natural storage
        substrate for a self-serve data mesh.
  \item \textbf{Federated computational governance}: global policies
        (GDPR compliance, CCPA data minimisation, HIPAA de-identification,
        data retention limits) are enforced \emph{automatically} at the
        platform level for all data products. Domain teams cannot opt out
        of global policies, but retain local flexibility in how they
        structure and process their data.
\end{enumerate}
\end{defbox}

\begin{caution}{Data Mesh Is Not a Technology}
Data Mesh is frequently misunderstood as a distributed storage architecture
or a multi-cloud strategy. It is neither. It is an organisational design
pattern for managing data responsibility. You can implement Data Mesh on
a single cloud with a single Lakehouse. Conversely, building a distributed
multi-cloud storage system does not give you a Data Mesh if the data
ownership and accountability structure is still centralised. The
transformation is primarily social (who is accountable for data quality?)
and secondarily technical (how does the infrastructure enforce that
accountability?).
\end{caution}

%%====================================================================
\section{MLOps: Operationalising Machine Learning}
\label{sec:mlops}
%%====================================================================

A trained machine learning model that exists only in a Jupyter notebook
delivers zero business value. The discipline of \term{MLOps} (Machine
Learning Operations) covers the engineering required to reliably deploy,
monitor, and retrain models in production. The first paper to formalise
this engineering gap was \autocite{MLTechnicalDebt2015}.

%--------------------------------------------------------------------
\subsection{Hidden Technical Debt in ML Systems}
\label{subsec:technical-debt}
%--------------------------------------------------------------------

Sculley et al.\ identified a fundamental structural problem with production
ML systems: the ML code itself---the model training loop---is only a tiny
fraction of the total codebase. The surrounding infrastructure is vast,
complex, and poorly understood.

\begin{keyinsight}{Hidden Technical Debt (Sculley et al., NeurIPS 2015)}
``Only a small fraction of real-world ML systems is composed of the ML
code. The required surrounding infrastructure is vast and complex.''
The paper identified nine categories of technical debt specific to ML
systems:
\begin{enumerate}
  \item \textbf{Entanglement}: changing one input feature changes the
        behaviour of all others (``Changing Anything Changes Everything'').
  \item \textbf{Hidden feedback loops}: a model's predictions influence
        the data it is next trained on---a source of instability invisible
        to offline evaluation.
  \item \textbf{Undeclared consumers}: other systems silently depend on
        a model's output format; schema changes break them invisibly.
  \item \textbf{Data dependency debt}: stale, underutilised, or
        epsilon-feature dependencies accumulate and become expensive to
        remove.
  \item \textbf{Configuration debt}: hyperparameters, thresholds, and
        pipeline configurations are not version-controlled.
  \item \textbf{Fixed thresholds in dynamic worlds}: a fraud threshold
        set in 2020 may be inappropriate in 2025 as transaction
        distributions shift.
  \item \textbf{Monitoring gaps}: the model trains on historical data
        but serves live data; distribution drift goes undetected until
        a business metric degrades.
  \item \textbf{Training--serving skew}: the features computed at training
        time use different code than the features computed at serving time,
        producing silent prediction errors.
  \item \textbf{Pipeline jungle}: ad-hoc data pipelines accumulate into
        an unmaintainable collection of scripts.
\end{enumerate}
MLOps is the systematic response to this debt: a set of tools and
practices that makes each failure mode visible, measurable, and preventable.
\end{keyinsight}

%--------------------------------------------------------------------
\subsection{The MLOps Pipeline}
\label{subsec:mlops-pipeline}
%--------------------------------------------------------------------

A complete MLOps pipeline has seven components, each addressing a specific
category of technical debt.

\textbf{1. Data ingestion} (Kafka + Spark): reproducible, versioned data
pipelines replace ad-hoc scripts. Data is written to the Lakehouse as a
Delta table, providing immutable, versioned input to all downstream stages.

\textbf{2. Feature engineering} (Apache Airflow DAGs): feature computation
logic is expressed as versioned, tested Airflow DAGs rather than one-off
notebook cells. A DAG that computes ``average transaction amount per user
over the past 30 days'' runs on the same code path in both the training
pipeline and the real-time serving pipeline, eliminating training--serving
skew.

\textbf{3. Feature store} (\tool{Feast}, \tool{Tecton}): a \term{feature
store} is a centralised repository that stores pre-computed features and
serves them to both training jobs (historical feature values for a training
window) and inference services (real-time feature values for a prediction
request). Without a feature store, different teams independently compute
the same features with subtly different logic, creating inconsistent model
inputs. The feature store also prevents \emph{point-in-time correctness}
violations: a training example for a customer's state at time $t$ must
use only feature values known at time $t$, not future values
(``data leakage'').

\textbf{4. Model training} (Spark MLlib, PyTorch, scikit-learn): the
model training code reads features from the feature store, trains the
model, and records all inputs (dataset version, feature set version,
hyperparameters) and outputs (model artefacts, evaluation metrics) in
an experiment tracker.

\textbf{5. Experiment tracking} (\tool{MLflow}): MLflow's \code{mlflow.log\_param()},
\code{mlflow.log\_metric()}, and \code{mlflow.log\_artifact()} API records
every experiment run's configuration and output. The MLflow Tracking UI
allows comparing runs by metric and reproducing any past result from
its logged configuration.

\textbf{6. Model registry} (MLflow Model Registry): trained models are
registered in a version-controlled catalogue with stage labels
(\texttt{Staging}, \texttt{Production}, \texttt{Archived}).
A model transitions from \texttt{Staging} to \texttt{Production} only
after passing automated integration tests and a manual approval step.
A failing production model can be rolled back to a previous version in
one command.

\textbf{7. Serving and monitoring} (\tool{KServe}, \tool{Seldon},
\tool{Evidently}): the model is deployed as a REST API endpoint with
canary deployment (10\% of traffic goes to the new model; 90\% stays
on the old model until the new model proves itself). \tool{Evidently}
or \tool{Arize} monitors the live feature distribution against the training
distribution, alerting when \term{data drift} is detected.

\begin{lstlisting}[style=pyspark, caption={MLflow experiment tracking and
  model registration in PySpark}]
import mlflow
import mlflow.spark
from pyspark.sql import SparkSession
from pyspark.ml.classification import GBTClassifier
from pyspark.ml.feature import VectorAssembler
from pyspark.ml.evaluation import BinaryClassificationEvaluator

spark = SparkSession.builder.appName("MLflowDemo").getOrCreate()
mlflow.set_tracking_uri("http://mlflow-server:5000")
mlflow.set_experiment("payment-fraud-detection-v2")

# Load feature-engineered training data from Delta Lake
train = spark.read.format("delta") \
             .load("hdfs:///lakehouse/features/fraud_train/") \
             .na.drop()

feature_cols = ["amount_usd", "txn_per_hour",
                "avg_txn_30d", "is_new_receiver",
                "hour_of_day", "day_of_week"]
assembler = VectorAssembler(inputCols=feature_cols,
                            outputCol="features")
train_vec = assembler.transform(train)

with mlflow.start_run(run_name="gbt_maxDepth5_lr01"):

    # --- Hyperparameters ---
    max_depth  = 5
    lr         = 0.1
    n_trees    = 100
    mlflow.log_param("maxDepth",         max_depth)
    mlflow.log_param("learningRate",     lr)
    mlflow.log_param("numTrees",         n_trees)
    mlflow.log_param("feature_set",      "v2.3")
    mlflow.log_param("training_version", 47)    # Delta Lake version

    # --- Train ---
    gbt = GBTClassifier(
        labelCol        = "is_fraud",
        featuresCol     = "features",
        maxDepth        = max_depth,
        stepSize        = lr,
        maxIter         = n_trees,
    )
    model = gbt.fit(train_vec)

    # --- Evaluate ---
    test    = assembler.transform(
        spark.read.format("delta")
             .load("hdfs:///lakehouse/features/fraud_test/").na.drop())
    preds   = model.transform(test)
    eval_   = BinaryClassificationEvaluator(labelCol="is_fraud")
    auroc   = eval_.evaluate(preds)
    mlflow.log_metric("auroc", auroc)
    print(f"AUROC = {auroc:.4f}")

    # --- Log model artefact ---
    mlflow.spark.log_model(model, "gbt_fraud_model",
                           registered_model_name="payment-fraud-gbt")

# --- Transition to Staging via Model Registry API ---
client = mlflow.MlflowClient()
latest = client.get_latest_versions("payment-fraud-gbt",
                                     stages=["None"])[0]
client.transition_model_version_stage(
    name    = "payment-fraud-gbt",
    version = latest.version,
    stage   = "Staging",
)
print(f"Model v{latest.version} promoted to Staging")
\end{lstlisting}

%--------------------------------------------------------------------
\subsection{Feature Stores and Training--Serving Skew}
\label{subsec:feature-store}
%--------------------------------------------------------------------

\term{Training--serving skew} is the most common silent failure mode in
production ML systems: the code that computes a feature at training time
differs from the code that computes the same feature at serving time,
because they were written by different teams at different times. The result
is a model that performs well on the training and test sets but degrades
mysteriously in production---because the features it sees during inference
are not the same features it was trained on.

A feature store eliminates this skew by maintaining a \emph{single}
implementation of each feature that is called by both the training pipeline
(to build the historical feature matrix) and the serving pipeline (to
compute real-time features for an incoming prediction request).

\begin{systemdesign}{Feature Store Architecture with Feast}
\tool{Feast} is an open-source feature store with two stores:
\begin{itemize}
  \item \textbf{Offline store} (Delta Lake / BigQuery / Redshift):
        stores the full historical record of each feature, indexed by
        \texttt{entity\_id} and \texttt{event\_timestamp}. Training
        jobs retrieve a point-in-time correct feature matrix by joining
        the entity's label events to the offline store at the exact
        timestamps of those events---preventing future data leakage.
  \item \textbf{Online store} (Redis / DynamoDB): stores only the
        \emph{latest} value of each feature for each entity. Inference
        requests look up the online store in $<$5 milliseconds. The
        offline store is periodically materialised to the online store
        (``feature materialisation'').
\end{itemize}
Both stores share the same feature transformation code, guaranteeing
that training and serving use identical feature values.
\end{systemdesign}

%%====================================================================
\section{Edge Analytics and Cloud-Native Big Data}
\label{sec:edge-cloud}
%%====================================================================

%--------------------------------------------------------------------
\subsection{The Three-Tier Edge--Cloud Architecture}
\label{subsec:three-tier}
%--------------------------------------------------------------------

IoT devices generate data at a rate that cannot economically be
transmitted in full to a central cloud cluster. A DESCO power grid
substation generates 100{,}000 voltage readings per second; a factory
floor with 500 vibration sensors generates 10 million samples per second.
Transmitting all this raw data to the cloud would require terabytes per
day of bandwidth---prohibitively expensive and subject to network latency
that makes real-time anomaly response impossible.

\term{Edge analytics} solves this by placing computation at three tiers.

\begin{defbox}{Three-Tier Edge--Cloud Architecture}
\begin{enumerate}
  \item \textbf{Tier 1: Device} (microcontroller, FPGA, SBC).
        At the data source, lightweight models (TensorFlow Lite, ONNX
        Runtime) perform preprocessing, filtering, and anomaly detection.
        Only \emph{events} (anomalies detected, thresholds exceeded) are
        forwarded upstream---not raw sensor readings. A cardiac monitor
        runs a QRS peak detection algorithm on-device and transmits only
        the detected heartbeat events, not the raw 500~Hz ECG waveform.
  \item \textbf{Tier 2: Edge Gateway} (local server, AWS Greengrass,
        Azure IoT Edge, NVIDIA Jetson).
        The edge gateway runs pre-trained inference models and windowed
        aggregation on the stream of device events from multiple devices.
        It forwards only \emph{summary alerts} or \emph{windowed
        aggregates} to the cloud. Typical data reduction ratio: $1000\times$
        (from raw 100{,}000 readings/sec to $\sim$100 alerts/sec).
        A factory's Greengrass gateway runs Isolation Forest on vibration
        sensor events to detect bearing wear; it sends one alert per
        anomalous 5-minute window rather than 30 million raw readings.
  \item \textbf{Tier 3: Cloud Lakehouse} (AWS S3 + Delta Lake,
        Azure ADLS + Iceberg).
        The cloud receives the aggregated alerts and periodically batched
        raw samples. It retrains global models on aggregated data from all
        edge nodes; long-term historical storage enables trend analysis
        and root-cause investigation. Updated models are pushed back
        down to the edge gateways.
\end{enumerate}
The upstream data flow is \textbf{Kafka} or \textbf{MQTT} between tiers.
The downstream model push is typically a CI/CD pipeline triggered by
MLflow model registry promotion.
\end{defbox}

\textbf{Cloud-native big data.} On-premise Hadoop clusters are being
replaced by \emph{managed cloud services} that decouple storage from
compute:
\begin{itemize}
  \item \textbf{Storage}: Amazon S3, Azure Data Lake Storage (ADLS),
        Google Cloud Storage (GCS) --- $\sim$\$23/TB-month, 11 nines
        durability, no cluster to maintain.
  \item \textbf{Compute}: AWS EMR (Elastic MapReduce), Google Dataproc,
        Azure HDInsight --- ephemeral Spark clusters that start in 5
        minutes, process the job, and terminate automatically (scaling
        to zero when idle). A batch pipeline that previously required
        a 100-node always-on cluster now runs on a 20-node cluster for
        3 hours per day---a 60--80\% cost reduction.
\end{itemize}

The Lakehouse is the natural convergence of cloud-native storage and
managed compute: Delta tables live on S3; ephemeral Spark clusters read
them directly, without copying data to a warehouse.

%%====================================================================
\section{Privacy-Preserving Analytics}
\label{sec:privacy}
%%====================================================================

Chapter~\ref{ch:domain-apps} introduced differential privacy and federated
learning in the healthcare context. This section formalises both concepts
and extends them to the general data engineering setting.

%--------------------------------------------------------------------
\subsection{Differential Privacy}
\label{subsec:dp}
%--------------------------------------------------------------------

\begin{defbox}{$(\varepsilon, \delta)$-Differential Privacy}
A randomised mechanism $\mathcal{M}: \mathcal{D} \to \mathcal{R}$ is
\textbf{$(\varepsilon, \delta)$-differentially private} if for all pairs
of neighbouring datasets $D, D'$ that differ in exactly one record, and
for all subsets of outputs $S \subseteq \mathcal{R}$:
\[
  \Pr[\mathcal{M}(D) \in S]
  \leq e^\varepsilon \cdot \Pr[\mathcal{M}(D') \in S] + \delta
\]
\textbf{Interpretation}: the probability that the mechanism's output
changes when one individual's record is added or removed is bounded by
$e^\varepsilon$ (multiplicatively) with probability $1 - \delta$.
Smaller $\varepsilon$ is stronger privacy. $\varepsilon = 0$ means
perfect privacy (the output reveals nothing about any individual);
$\varepsilon = \infty$ means no privacy guarantee.

\textbf{The Laplace mechanism}: for a real-valued function $f: \mathcal{D}
\to \mathbb{R}$ with sensitivity $\Delta f = \max_{D, D'} |f(D) - f(D')|$,
adding Laplace noise achieves $(\varepsilon, 0)$-differential privacy:
\[
  \mathcal{M}(D) = f(D) + \text{Lap}\!\left(\frac{\Delta f}{\varepsilon}\right)
\]
\textbf{The Gaussian mechanism}: adding Gaussian noise achieves
$(\varepsilon, \delta)$-differential privacy, which is practically
preferable for machine learning (gradient perturbation).
\end{defbox}

\begin{lstlisting}[style=pyspark, caption={Differential privacy with the
  Laplace mechanism: publishing district-level case counts}]
import numpy as np

def laplace_mechanism(true_count: int,
                      sensitivity: float,
                      epsilon: float,
                      rng: np.random.Generator) -> float:
    """Add Laplace noise to a count query."""
    scale = sensitivity / epsilon
    noise = rng.laplace(loc=0.0, scale=scale)
    # Counts must be non-negative; clip to 0
    return max(0.0, true_count + noise)

rng     = np.random.default_rng(2024)
epsilon = 1.0      # privacy budget
sens    = 1.0      # sensitivity of a COUNT query

# True dengue case counts by region (confidential)
true_counts = {
    "New York":    1_402,
    "London":      387,
    "Berlin":      201,
    "Sydney":      94,
    "Toronto":     47,
}

print(f"epsilon = {epsilon} (stronger privacy -> larger noise)")
print(f"{'Region':<14} {'True':>8} {'Published':>12} {'Error':>8}")
print("-" * 46)

for region, count in true_counts.items():
    published = laplace_mechanism(count, sens, epsilon, rng)
    error     = abs(published - count)
    print(f"{region:<14} {count:>8} {published:>12.1f} {error:>8.1f}")

# Compare: tighter budget -> more noise
for eps in [0.1, 0.5, 1.0, 2.0, 5.0]:
    noised = laplace_mechanism(1_402, 1.0, eps, rng)
    print(f"  epsilon={eps:<4}: published New York count = {noised:.1f}")
\end{lstlisting}

\textbf{Real-world deployments.}
\begin{itemize}
  \item \textbf{Apple}: uses local differential privacy (LDP) for
        keyboard usage statistics on iOS---noise is added on the
        device before any data leaves the phone.
  \item \textbf{Google}: RAPPOR (Randomised Aggregatable Privacy-Preserving
        Ordinal Response) for Chrome usage telemetry.
  \item \textbf{US Census Bureau}: the 2020 US Census used differential
        privacy (the TopDown algorithm) to prevent re-identification of
        individuals in block-level population tables.
\end{itemize}

%--------------------------------------------------------------------
\subsection{Federated Learning}
\label{subsec:federated-learning}
%--------------------------------------------------------------------

\autocite{McMahan2017} introduced \term{federated learning} (FL) as a
method for training a shared model on data held across multiple devices
or institutions without centralising the raw data. The core algorithm,
\emph{FedAvg}, works as follows:
\begin{enumerate}
  \item A \textbf{central server} initialises a global model $w_0$ and
        distributes it to all participating nodes (hospitals, mobile devices,
        banks).
  \item Each node trains the model locally for $E$ epochs on its own data,
        producing a locally updated model $w_i^{(t)}$.
  \item Each node sends only its \emph{model weight update}
        $\Delta w_i = w_i^{(t)} - w_0^{(t)}$---not its raw data---to
        the server.
  \item The server aggregates the updates using weighted averaging:
        \[
          w_0^{(t+1)} = \sum_{i=1}^{K} \frac{n_i}{N} \Delta w_i
        \]
        where $n_i$ is node $i$'s dataset size and $N = \sum n_i$.
  \item Repeat for $T$ global rounds.
\end{enumerate}

The raw data never leaves the originating node. This enables model training
across organisational boundaries where data sharing is legally prohibited.

\begin{techdebt}{Federated Learning Limitations in Production}
Despite its privacy advantages, federated learning introduces engineering
challenges that must be addressed before production deployment:
\begin{enumerate}
  \item \textbf{Statistical heterogeneity (non-IID data)}: each node's
        data reflects local conditions (a rural hospital's patient population
        differs from an urban hospital's). Naive FedAvg converges slowly
        or to a poor global model when data distributions are highly
        non-IID. Techniques such as FedProx (adding a proximal term to
        the local loss) and SCAFFOLD (correction of client drift) address
        this.
  \item \textbf{Communication cost}: transmitting full model weight vectors
        (sometimes hundreds of millions of parameters) after every round
        is expensive. Gradient compression, quantisation, and sparsification
        reduce communication by $10$--$100\times$ at some accuracy cost.
  \item \textbf{Gradient inversion attacks}: recent research has shown
        that model weight updates can leak information about the training
        data. Combining federated learning with differential privacy
        (DP-SGD: adding noise to gradients before uploading) provides
        formal privacy guarantees even against a malicious server.
  \item \textbf{Partial participation}: in cross-device federated learning
        (e.g., mobile phones), only a fraction of nodes are available at
        each round (device is off, battery is low). The training algorithm
        must be robust to arbitrary dropouts.
\end{enumerate}
\end{techdebt}

%%====================================================================
\section{The Dataflow Model and Apache Beam}
\label{sec:dataflow}
%%====================================================================

Chapter~\ref{ch:pipelines} described the Lambda architecture as the
practitioner's solution to the batch--stream unification problem: run a
batch pipeline for accuracy and a stream pipeline for latency, then merge
their outputs in a serving layer. The Lambda architecture's fundamental
flaw is its \emph{dual code path}: identical business logic must be
maintained in two separate codebases that gradually drift.

\autocite{Dataflow2015} proposed a more principled solution: a single
programming model that is expressive enough to describe both batch and
streaming workloads, treating batch as a special case of streaming.

%--------------------------------------------------------------------
\subsection{The Four Questions}
\label{subsec:four-questions}
%--------------------------------------------------------------------

Any data pipeline can be fully characterised by four questions:

\begin{defbox}{The Dataflow Model's Four Ws (Akidau et al., VLDB 2015)}
\begin{enumerate}
  \item \textbf{What} results are computed?
        The \emph{transformation} applied to the data: sum, count,
        average, join, custom function. This is equivalent to the
        \texttt{SELECT} clause in SQL.
  \item \textbf{Where} in event time are results computed?
        The \emph{windowing strategy} partitions the unbounded event
        stream into finite chunks:
        \begin{itemize}
          \item \textbf{Tumbling window}: fixed, non-overlapping windows
                (e.g., count transactions every 5 minutes).
          \item \textbf{Sliding window}: fixed-size, overlapping windows
                (e.g., 1-hour window updated every 5 minutes).
          \item \textbf{Session window}: variable-size window defined by
                activity---a window closes when there has been no event
                for a specified gap duration.
        \end{itemize}
  \item \textbf{When} in processing time are results materialised?
        The \emph{trigger} controls when the system emits a result for
        a window:
        \begin{itemize}
          \item \textbf{Watermark trigger}: emit when the watermark
                (the system's estimate of the latest event time it has
                seen minus an allowance for late arrivals) passes the
                window's end time.
          \item \textbf{Periodic trigger}: emit every $n$ seconds
                regardless of the watermark.
          \item \textbf{Count trigger}: emit after $n$ events.
          \item \textbf{Composite}: e.g., emit early every 10 seconds
                AND emit a final result when the watermark passes.
        \end{itemize}
  \item \textbf{How} do refinements relate to previous results?
        The \emph{accumulation mode} determines what happens when a
        window fires multiple times (early trigger, then final):
        \begin{itemize}
          \item \textbf{Discarding}: each firing emits only the new data
                arrived since the last firing.
          \item \textbf{Accumulating}: each firing emits the full running
                total since the window opened (overwrites previous output).
          \item \textbf{Accumulating and retracting}: each firing emits
                both the new value AND a retraction of the previous value
                (for downstream systems that need to subtract the old value).
        \end{itemize}
\end{enumerate}
\textbf{Batch as a special case of streaming}: in this model, a batch
job is a streaming job with a single tumbling window covering the entire
dataset, a watermark trigger set at the end of the input, and discarding
accumulation. The same code expresses both.
\end{defbox}

%--------------------------------------------------------------------
\subsection{Apache Beam: One Pipeline, Multiple Runners}
\label{subsec:beam}
%--------------------------------------------------------------------

\tool{Apache Beam} (2016) is the open-source implementation of the
Dataflow Model. A Beam pipeline is written in Python, Java, or Go using
the Beam SDK. The same pipeline code then runs on multiple \emph{runners}:
\begin{itemize}
  \item \textbf{Apache Spark} runner: for batch processing on HDFS or S3.
  \item \textbf{Apache Flink} runner: for low-latency streaming.
  \item \textbf{Google Cloud Dataflow}: fully managed, auto-scaling stream
        and batch processing on GCP.
  \item \textbf{Direct runner}: for local testing without a cluster.
\end{itemize}

\begin{lstlisting}[style=pyspark, caption={Apache Beam: fraud detection
  pipeline with tumbling windows (same code runs on Spark or Flink)}]
import apache_beam as beam
from apache_beam.options.pipeline_options import PipelineOptions
from apache_beam.transforms.window import FixedWindows
from apache_beam.transforms.trigger import (AfterWatermark,
    AfterProcessingTime, AccumulationMode)
import json

# --- Configuration: change runner to switch between Spark / Flink / local ---
options = PipelineOptions([
    "--runner=FlinkRunner",
    "--flink_master=flink-master:6123",
    # Use SparkRunner for batch:
    # "--runner=SparkRunner", "--spark_master=spark://...",
])

def parse_event(raw: bytes) -> dict:
    """Deserialise a Kafka event bytes to dict."""
    return json.loads(raw.decode("utf-8"))

def extract_key_amount(record: dict):
    """Emit (region, amount_usd) pairs."""
    return (record.get("region", "Unknown"),
            float(record.get("amount_usd", 0.0)))

class SumValues(beam.CombineFn):
    def create_accumulator(self): return 0.0
    def add_input(self, acc, element): return acc + element
    def merge_accumulators(self, accs): return sum(accs)
    def extract_output(self, acc): return acc

def format_output(kv) -> str:
    region, total = kv
    return f"{region}: USD {total:,.0f}"

with beam.Pipeline(options=options) as p:
    result = (
        p
        # WHAT: read from Kafka
        | "ReadKafka" >> beam.io.ReadFromKafka(
            consumer_config={
                "bootstrap.servers": "kafka:9092",
                "group.id": "beam-fraud-pipeline",
            },
            topics=["payment-events"],
        )
        | "ParseJSON"      >> beam.Map(parse_event)
        | "ExtractPairs"   >> beam.Map(extract_key_amount)

        # WHERE: 5-minute tumbling event-time windows
        | "Window" >> beam.WindowInto(
            FixedWindows(5 * 60),        # 5-minute windows
            trigger=AfterWatermark(
                early=AfterProcessingTime(30),   # emit every 30s early
            ),
            accumulation_mode=AccumulationMode.ACCUMULATING,
            allowed_lateness=2 * 60,     # accept events up to 2 min late
        )

        # WHAT: aggregate amount per region per window
        | "SumByRegion"    >> beam.CombinePerKey(SumValues())
        | "Format"         >> beam.Map(format_output)

        # Output: write aggregates to the Lakehouse (or console for testing)
        | "Write" >> beam.io.WriteToText("hdfs:///lakehouse/fraud/agg")
    )
\end{lstlisting}

%%====================================================================
\section{Research Spotlight}
\label{sec:ch11-research}
%%====================================================================

\begin{researchspotlight}{Armbrust et al.\ (2020) · Armbrust, Ghodsi, Xin
  \& Zaharia (2021) · Akidau et al.\ (2015) · Sculley et al.\ (2015) ---
  The Modern Data Engineering Stack}

\textbf{Work 1: Armbrust, M., et al.\ (2020). ``Delta Lake: High-Performance
ACID Table Storage over Cloud Object Stores.'' \textit{Proceedings of the
VLDB Endowment}, 13(12), 3411--3424.} (Databricks / Apache Foundation;
Rank A* venue.)

Delta Lake addresses a fundamental limitation of cloud object stores
(S3, Azure ADLS, GCS): they provide no ACID semantics. A concurrent
write from two Spark executors to the same S3 prefix produces undefined
results; a partial write due to a job failure leaves the table in a
corrupt state; there is no mechanism for rolling back a bad commit.

The paper's solution is the \emph{Delta Log}: an ordered sequence of
atomic JSON commits stored inside the table directory in S3. Each commit
records added and removed Parquet files, schema changes, and metadata.
Readers reconstruct the table state by reading the log; they always see
a consistent snapshot. Writers implement optimistic concurrency control
by racing to write the next log entry number---only one succeeds.
Time travel is a consequence of the log's durability: any past version
is recoverable by reading the log up to that version's commit.

Three benchmarks from the paper:
\begin{enumerate}
  \item Write throughput equal to raw Parquet on S3 (the log adds $<$2\%
        overhead for writes; checkpoint files amortise read cost).
  \item Z-ordering reduces the number of files scanned by queries with
        two-column filters by 95\% compared to Hive partitioning on
        a single column.
  \item MERGE INTO (upsert) is $10\times$ faster than the equivalent
        read--filter--rewrite pattern on raw Parquet.
\end{enumerate}

\textbf{Impact}: Delta Lake is now the default table format for Databricks
Lakehouse Platform, used by over 7{,}000 organisations. It has been
adopted as a core component of AWS Lake Formation and Azure Synapse
Analytics. Apache Iceberg (Netflix) and Apache Hudi (Uber) adopt the
same transaction-log principle with different trade-offs.

\medskip
\textbf{Work 2: Armbrust, M., Ghodsi, A., Xin, R., \& Zaharia, M. (2021).
``Lakehouse: A New Generation of Open Platforms That Unify Data
Warehousing and Advanced Analytics.'' \textit{11th Annual Conference on
Innovative Data Systems Research (CIDR 2021)}.} (Databricks; invited
vision paper.)

The Lakehouse paper formalises the observation that Delta Lake's ACID
semantics make it possible to use a data lake as a data warehouse---without
a separate copy of the data in a proprietary warehouse system.
Zaharia et al.\ benchmark the Lakehouse architecture (Delta Lake + Spark
SQL on S3) against the leading cloud data warehouses (Redshift, Snowflake,
BigQuery) on the TPC-DS benchmark:
\begin{itemize}
  \item At the 3~TB scale factor, Databricks SQL on Delta Lake achieves
        comparable performance to Redshift RA3 at 3--5$\times$ lower cost.
  \item The ``no ETL copy'' property means analysts query data that is
        minutes old, not 12 hours old as in the traditional lake-to-warehouse
        pipeline.
\end{itemize}
The paper formally defines the Lakehouse using five properties (open
storage, metadata layer, SQL performance, ML support, governance)
and identifies DuckDB as an in-process analytical engine that further
democratises Lakehouse query access without a Spark cluster.

\textbf{Impact}: Snowflake, Google, and Microsoft have all published
blog posts and white papers in response to the Lakehouse paper.
Snowflake's Iceberg integration, Google's BigLake, and AWS Lake Formation
are all responses to the competitive pressure created by this paper.

\medskip
\textbf{Work 3: Akidau, T., et al.\ (2015). ``The Dataflow Model: A
Practical Approach to Balancing Correctness, Latency, and Cost in
Massive-Scale, Unbounded, Out-of-Order Data Processing.''
\textit{Proceedings of the VLDB Endowment}, 8(12), 1792--1803.}
(Google; VLDB Best Paper Award.)

Google's internal production streaming system (MillWheel, 2013) revealed
two fundamental inadequacies of the Lambda architecture: maintaining dual
code paths introduced synchronisation bugs; and the existing streaming
model (Storm, early Spark Streaming) provided no principled way to handle
out-of-order events---events that arrive at the processing system after
the window covering their event time has already been emitted.

Akidau et al.'s Dataflow Model answers four questions for any data pipeline
(What, Where, When, How) and introduces the \emph{watermark} as a
first-class primitive: the watermark is the system's estimate of the
maximum event time for which all earlier events have been received.
Triggers allow windows to emit early results (for latency) while
waiting for late arrivals (for completeness). The accumulation mode
controls how late-arriving events update previously emitted results.

\textbf{Impact}: The Dataflow Model was open-sourced as Apache Beam in
2016. Apache Flink adopted the Dataflow Model's event-time windowing
semantics directly (the \texttt{EventTimeSessionWindows} and
\texttt{EventTimeTrigger} APIs are direct implementations of the Four
Ws). Google Cloud Dataflow is the managed production deployment.
Spark Structured Streaming's \texttt{watermark} API
(Chapter~\ref{ch:pipelines}) is an implementation of the same model.

\medskip
\textbf{Work 4: Sculley, D., et al.\ (2015). ``Hidden Technical Debt in
Machine Learning Systems.'' \textit{Advances in Neural Information
Processing Systems (NeurIPS 2015)}.} ($>$5{,}000 citations; Google.)

Written by Google engineers managing internal ML production systems,
this paper is one of the most practically influential papers of the
2010s. Its central argument is that the visible ML code---the model
training loop---is only a small fraction of the actual cost of a
production ML system. Nine categories of hidden technical debt
(entanglement, undeclared consumers, data dependency debt, configuration
debt, fixed thresholds, monitoring gaps, training--serving skew,
pipeline jungle, world changes) account for far more engineering
effort than the ML code itself.

The paper's impact was to crystallise the \emph{MLOps} discipline:
toolmakers (Databricks MLflow, Tecton, Weights \& Biases, Evidently AI)
built their entire product around eliminating specific categories of
the debt described in this paper. Stanford's CS329S (ML Systems) assigns
the paper as the first required reading and maps its debt taxonomy to
every tool in the MLOps stack.
\end{researchspotlight}

%%====================================================================
\section{Case Study: Enterprise Lakehouse Migration in Production}
\label{sec:lakehouse-case}
%%====================================================================

\begin{casestudy}{From Lambda Architecture to Lakehouse ---
  Pinterest, Comcast, and the Production Engineering Lessons}

\textbf{Context: Lambda Pain at Scale.}
By 2018--2020, organisations running the Lambda architecture at petabyte
scale were experiencing the same compounding pain:
\begin{itemize}
  \item \textbf{Pinterest (300+ PB stored, 10+ billion events/day)}: for
        every business metric (engagement rate, recommendation click-through,
        ad impression count), two separate codebases existed---a Spark batch
        job for the weekly reports and a Flink streaming job for the
        real-time dashboard. Engineers estimated they spent 40\% of their
        time keeping the two code paths in sync. When business logic changed
        (new definition of an ``engagement event''), both codebases had to
        be updated simultaneously---and failures in synchronisation produced
        inconsistent numbers that took days to diagnose.
  \item \textbf{Comcast (100 TB of new data per day)}: data was ingested
        into Amazon S3 (the data lake) via Spark batch jobs. Every night, a
        separate ETL pipeline extracted the previous day's data from S3 and
        loaded it into Amazon Redshift (the data warehouse). Business analysts
        queried Redshift; ML teams queried S3. The copy introduced a 12-hour
        data freshness lag for analysts, and the two stores inevitably diverged
        as some records were cleaned in the ETL pipeline and others were not.
\end{itemize}

\textbf{Migration Architecture.}
Both organisations migrated to a Lakehouse using Delta Lake on Amazon S3 +
Databricks (Spark) as the unified query engine. The migration strategy
was \emph{table-by-table}, not a big-bang cutover:
\begin{enumerate}
  \item Convert the highest-value, most-queried table first (e.g., the
        daily active users table). Enable schema enforcement immediately.
        Run the Delta version and the legacy version in parallel for two
        weeks, comparing query results.
  \item Migrate the Flink streaming pipeline to write to Delta using
        Spark Structured Streaming with the Delta Lake sink. The streaming
        table and the batch table are now the same Delta table---the same
        Parquet files, the same ACID log.
  \item Redirect BI tools (Tableau, Superset) from Redshift to Delta Lake
        via Spark SQL. Validate dashboard numbers against the Redshift
        baseline.
  \item Decommission the Redshift cluster and the nightly ETL pipeline.
\end{enumerate}

\textbf{Measured Outcomes.}
\begin{center}
\small
\renewcommand{\arraystretch}{1.25}
\begin{tabular}{lll}
  \toprule
  \textbf{Metric} & \textbf{Before (Lambda)} & \textbf{After (Lakehouse)} \\
  \midrule
  ETL pipeline count per metric  & 3--5 jobs   & 1 job \\
  BI data freshness              & 12 hours    & $<$15 minutes \\
  Storage cost per TB-month      & \$23 lake + \$230 DWH & \$23 (lake only) \\
  ML training data access        & File copy from S3 & Direct Delta read \\
  Schema violations detected     & After landing & Blocked at ingest \\
  GDPR record deletion           & Rewrite all Parquet files & Single delta log entry \\
  Analyst query latency (P95)    & 8 minutes (Hive) & 45 seconds (Delta) \\
  Code paths per pipeline        & 2 (batch + stream) & 1 \\
  \bottomrule
\end{tabular}
\end{center}

\textbf{Five Engineering Lessons.}
\begin{enumerate}
  \item \textbf{Enable schema enforcement on day one.}
        Teams that enabled schema enforcement in the first week of migration
        discovered and fixed 90\% of their data quality bugs within that
        week. Teams that deferred schema enforcement ``until the migration
        was stable'' spent months debugging corrupted tables that had
        accumulated malformed records during the deferral period.
  \item \textbf{Time travel replaced manual snapshot exports.}
        Compliance teams had previously maintained a weekly manual export
        of the Redshift table to a separate S3 folder as an ``audit snapshot.''
        After migration, \texttt{VERSION AS OF} or \texttt{TIMESTAMP AS OF}
        queries served this function automatically---eliminating 6~TB/month
        of redundant storage.
  \item \textbf{Z-ordering beats blanket partitioning.}
        The legacy Hive tables were partitioned by \texttt{date}---reasonable
        for time-range queries but inadequate for the common multi-column
        filter \texttt{WHERE user\_id = X AND date BETWEEN Y AND Z}
        (which still scanned the entire date partition). Z-ordering on
        \texttt{(user\_id, date)} reduced scanned files by 78--85\% for
        these queries.
  \item \textbf{Migration, not replacement.}
        The table-by-table migration strategy eliminated downtime and allowed
        the team to validate each table independently. A big-bang migration
        that switches all tables simultaneously produces a situation where
        a bug in one table poisons all downstream consumers and is impossible
        to roll back cleanly.
  \item \textbf{The bottleneck shifted from storage to governance.}
        After the technical migration, the dominant remaining challenge was
        \emph{organisational}: who owns which Delta table? Which teams can
        write to it? Who is responsible when a table's data quality degrades?
        These questions---the subject of the Data Mesh paradigm---proved
        harder to solve than the technical migration itself.
\end{enumerate}
\end{casestudy}

%%====================================================================
\section*{Chapter Summary}
\label{sec:ch11-summary}
%%====================================================================

This chapter described the fourth era of big data technology and the
architecture that defines it.

The Lakehouse combines cheap open-format object storage with an open
table format (Delta Lake, Apache Iceberg, or Apache Hudi) that provides
ACID transactions via a transaction log, time travel via log durability,
and schema enforcement at the storage layer. Delta Lake's optimistic
concurrency control achieves serialisable writes on eventually-consistent
S3 without a dedicated lock manager. The Lakehouse eliminates the
lake--warehouse split: BI and ML workloads read the same ACID-consistent
data, and the Lambda architecture's 12-hour freshness lag disappears.

The Data Mesh addresses the organisational bottleneck created when a
central platform team owns all data pipelines: it decentralises ownership
to domain teams (domain ownership), standardises the interface of each
dataset (data as a product), provides infrastructure abstractions that
non-engineers can use (self-serve platform), and enforces global
compliance policies automatically (federated governance).

MLOps makes production ML reliable by addressing the hidden technical
debt identified by Sculley et al.: feature stores eliminate training--serving
skew; MLflow tracks experiment configurations and enables result reproduction;
model registries provide versioned, staged deployment; and monitoring detects
data drift before it degrades business outcomes.

Edge analytics extends distributed processing to three tiers: lightweight
on-device inference (TensorFlow Lite, ONNX), edge gateway aggregation
(Greengrass, Azure IoT Edge), and cloud Lakehouse storage and model
retraining. Cloud-native managed services (AWS EMR, Google Dataproc)
decouple storage (S3, GCS) from compute, reducing cost by 60--80\%
compared to always-on HDFS clusters.

The Dataflow Model (Akidau et al., 2015) eliminates the Lambda
architecture's dual code path by providing a single programming model
with four dimensions: What (transformation), Where (windowing),
When (triggering), and How (accumulation). Apache Beam implements this
model and runs on Spark, Flink, and Google Cloud Dataflow---one pipeline
code, multiple backends.

%%====================================================================
\section*{Exercises}
\label{sec:ch11-exercises}
%%====================================================================

\exercisesection{Short Questions}

\sq Describe the four eras of big data technology. For each era, state
the primary engineering bottleneck that the previous era's tools could
not address and the key innovation that resolved it.

\sq What is the Delta Log in Delta Lake? How does it achieve atomicity,
isolation, and durability on Amazon S3, which provides no native locking
or atomic rename operations?

\sq A Spark job writes to a Delta table at the same time as a second Spark
job. Only one write succeeds; the other throws a \texttt{ConcurrentModificationException}.
Explain the optimistic concurrency control mechanism that produced this outcome.

\sq What is time travel in Delta Lake? Give two specific production use
cases where the ability to query a past version of a Delta table provides
direct business value.

\sq Describe the difference between Delta Lake's \textbf{schema enforcement}
and \textbf{schema evolution}. What problem does each solve, and what
is the consequence of not having schema enforcement in a data lake?

\sq What is Z-ordering? Explain why it outperforms single-column Hive
partitioning for queries that filter on two columns simultaneously.

\sq State the four principles of the Data Mesh paradigm. For each, identify
the specific bottleneck of centralised data architecture that it addresses.

\sq What is training--serving skew? Give a specific example from a fraud
detection context, describe how it produces silent prediction errors,
and explain how a feature store prevents it.

\sq Define $(\varepsilon, \delta)$-differential privacy formally. If
$\varepsilon = 0.1$ vs.\ $\varepsilon = 5.0$: which provides stronger
privacy, which adds more noise, and which is more appropriate for a
published aggregate statistic vs.\ an internal model training signal?

\sq Describe the Dataflow Model's four Ws. Explain how a batch MapReduce
job is a special case of the streaming model under this framework.

\sq What is a watermark in the Dataflow Model? Why is a watermark necessary
for correctly processing out-of-order streaming events?

\sq What is the ``no ETL copy'' advantage of the Lakehouse? Describe
the specific data freshness and cost problems it eliminates.

\exercisesection{Analytical Problems}

\ap \textbf{Delta Lake Concurrency Analysis.}
A payment transaction table is stored as a Delta table on S3. Three
concurrent Spark jobs attempt to write to the table simultaneously:
\begin{itemize}
  \item Job A: append 10 million new transactions (started first).
  \item Job B: UPDATE 500 transactions where \texttt{region='Northeast'}
        (started 2 seconds after Job A).
  \item Job C: DELETE all transactions where \texttt{amount\_usd < 0.50}
        (started 3 seconds after Job A).
\end{itemize}
\begin{enumerate}[label=(\alph*)]
  \item Which job writes the Delta Log entry numbered 00042.json if
        the current latest version is 41? Which jobs conflict?
  \item Under Delta Lake's OCC: if Job A succeeds, does Job B automatically
        fail, or does it check whether its operation conflicts with Job A's
        specific changes? Explain the conflict detection logic.
  \item If Job C fails with a \texttt{ConcurrentModificationException},
        what exactly does it retry? Does it re-read the data from scratch
        or simply retry the log append?
  \item After all three operations complete, how does a time-travel query
        \texttt{VERSION AS OF 41} behave? What data does it return?
\end{enumerate}

\ap \textbf{MLOps Debt Taxonomy.}
A hospital data science team built a sepsis prediction model in 2022.
It is deployed in production at a large academic medical centre and makes
predictions every 15 minutes for each ICU patient. In 2025, the team notices the model's
clinical alerts are being ignored by nurses, who say ``the model keeps
false-alarming.'' An investigation reveals:
\begin{itemize}
  \item The model was trained on 2019--2021 ICU data. In 2023, the hospital
        switched from manual vital-sign charting to an automated monitoring
        system that records vitals every 1 minute instead of every 4 hours.
        The model was not retrained.
  \item The real-time feature computation at serving time uses a slightly
        different rolling-window logic than the training pipeline (off-by-one
        in window boundary calculation).
  \item No monitoring was configured to detect feature distribution drift.
\end{itemize}
For each issue: (a)~identify the Sculley et al.\ technical debt category;
(b)~describe the specific MLOps tool or practice that would have prevented it;
(c)~estimate the severity of the clinical impact if the model's false positive
rate has risen from 12\% to 45\%.

\ap \textbf{Dataflow Model Design.}
Design a Beam pipeline for a real-time ride-sharing driver location analytics
system:
\begin{itemize}
  \item Source: GPS ping events arrive at Kafka at 10 events/second
        per active driver (estimated 50{,}000 active drivers during
        peak hour $\Rightarrow$ 500{,}000 events/second).
  \item Events have a 5--30 second event-time lag (mobile network latency).
  \item Required output: count of active drivers per city zone per
        5-minute window, emitted every 30 seconds with an early result,
        and a final result when the watermark passes the window end.
\end{itemize}
Answer the Four Ws:
\begin{enumerate}[label=(\alph*)]
  \item \textbf{What}: specify the transformation. What is the key,
        what is aggregated, and what is the output schema?
  \item \textbf{Where}: specify the windowing strategy. Why is a tumbling
        window preferable to a sliding window for this use case?
  \item \textbf{When}: specify the trigger. Write the trigger configuration
        in pseudocode (AfterWatermark / AfterProcessingTime / composite).
        What allowed lateness should be set given the 5--30s event-time lag?
  \item \textbf{How}: specify the accumulation mode. If the 30-second
        early trigger emits ``Downtown: 12{,}400 drivers'' and the final
        watermark trigger emits ``Downtown: 12{,}800 drivers'', what does the
        downstream system receive under \emph{Accumulating} mode vs.\
        \emph{Accumulating and Retracting} mode?
\end{enumerate}

\ap \textbf{Differential Privacy Budget Allocation.}
A public health agency plans to publish monthly region-level disease
statistics (dengue case counts, influenza hospitalisations, COVID-19 test
positivity rate) for 50 regions. Each statistic is a COUNT or AVERAGE
query over the patient records database. The total monthly privacy budget
is $\varepsilon_{\text{total}} = 2.0$.
\begin{enumerate}[label=(\alph*)]
  \item The agency wants to publish 5 statistics per region per month
        (dengue count, influenza count, COVID count, COVID positivity
        rate, and RSV count). How should the budget be allocated across
        these 5 queries using \emph{sequential composition} (each query
        spends from the budget independently)?
  \item A statistician proposes using \emph{parallel composition} instead:
        since the region-level counts are disjoint subsets of the total
        population, the budget does not need to be divided across regions.
        Explain why this is correct and what it implies for the per-region
        noise level.
  \item With $\varepsilon = 0.4$ per query (from part (a)), and a COUNT
        query sensitivity of $\Delta f = 1$, compute the standard deviation
        of the Laplace noise. Is this noise level acceptable for a region
        with 1{,}400 dengue cases? For a region with 14 cases?
        Discuss the small-count problem.
  \item Propose a solution for the small-count problem that maintains
        differential privacy while still publishing useful statistics for
        low-incidence regions.
\end{enumerate}

\exercisesection{Design Problems}

\dpr \textbf{National Mobile Payment Lakehouse (capstone design).}
Design a complete Lakehouse-based data analytics platform for a nationwide
mobile payment company (M-Pesa / PayPal scale: 50 million transactions/day,
5 years of historical data, 30 million registered users).

\begin{enumerate}[label=(\alph*)]
  \item \textbf{Ingestion layer}: specify the Kafka topic structure
        (topic, partition count, key, value schema), the Spark Structured
        Streaming sink configuration for writing to Delta Lake, and the
        trigger interval.
  \item \textbf{Storage design}: design the Delta table schema for the
        raw transaction table. Specify partitioning columns and Z-ordering
        columns. Justify each choice with reference to the query patterns
        (time-range queries for compliance, user-level queries for fraud,
        region-level aggregations for dashboards).
  \item \textbf{Analytical layers}: describe the curated table hierarchy
        (bronze $\to$ silver $\to$ gold) using Delta Lake. What
        transformations happen at each layer? What SLA (freshness,
        accuracy) applies to each layer?
  \item \textbf{MLOps}: design the fraud detection MLOps pipeline.
        Specify: feature store features (at least 5), MLflow experiment
        tracking configuration, model registry stages, serving endpoint
        latency target, and drift monitoring metric.
  \item \textbf{Governance}: specify the Data Mesh domain ownership
        structure for three domains (Payments, Compliance, Marketing).
        Which Delta tables does each domain own? What global governance
        policies are enforced automatically?
\end{enumerate}

\dpr \textbf{National Health Analytics Lakehouse.}
Design a Lakehouse for a national health analytics platform (e.g., NHS
England or US Health and Human Services), integrating data from 300
hospitals (HL7 FHIR), pharmacy chains, and community health workers
(mobile app events via Kafka).

\begin{enumerate}[label=(\alph*)]
  \item \textbf{Ingestion}: specify the Kafka Connect configuration for
        ingesting FHIR R4 Patient and Observation resources from 300
        hospitals. How do you handle hospitals with different FHIR
        versions (R3 vs R4)?
  \item \textbf{Privacy}: specify the de-identification pipeline that
        must run before any data enters the Lakehouse gold layer.
        Which HIPAA safe harbour identifiers are removed, and which are
        generalised (e.g., date-of-birth $\to$ age-band)?
  \item \textbf{Federated learning}: the health authority wants to train
        a sepsis prediction model without ever centralising individual
        patient records from private hospitals. Design the federated
        learning architecture: number of rounds, local epochs per round,
        FedAvg aggregation, and validation strategy.
  \item \textbf{Real-time alerting}: specify the Dataflow pipeline for
        detecting an influenza outbreak: window type and size, trigger,
        threshold (what counts as an alert), and output destination.
\end{enumerate}

\exercisesection{Programming Exercises}

\pe \textbf{Delta Lake ACID Operations.}
Using PySpark and the \texttt{delta-spark} library locally, implement
and verify Delta Lake's ACID properties:
\begin{enumerate}[label=(\alph*)]
  \item Create a Delta table for hospital admissions:
        \texttt{(patient\_id STRING, region STRING, diagnosis\_code STRING,
        admit\_date DATE, los\_days INT)}.
        Write 10{,}000 synthetic rows.
  \item Append 2{,}000 new records (Batch 2).
  \item Execute an UPDATE to add 2 to \texttt{los\_days} for all patients
        where \texttt{region = 'Northeast'}.
  \item Execute a DELETE for all patients where \texttt{los\_days < 1}
        (invalid records).
  \item Query \texttt{VERSION AS OF 1} (after Batch 2 but before the
        UPDATE) and count the Northeast patients. Verify the count matches
        what you expect.
  \item Run \texttt{DESCRIBE HISTORY} and display all four transaction
        entries. For each, identify the operation type.
  \item Run \texttt{OPTIMIZE} with Z-ordering on \texttt{(region,
        admit\_date)} and observe the reduction in the number of Parquet
        files.
\end{enumerate}

\pe \textbf{MLflow Experiment Tracking and Model Registry.}
Build a complete MLflow experiment for payment fraud classification:
\begin{enumerate}[label=(\alph*)]
  \item Generate 100{,}000 synthetic transaction records with
        \texttt{(amount\_usd, txn\_per\_hour, is\_new\_receiver,
        hour\_of\_day, is\_fraud)}.
  \item Train three models (Logistic Regression, Random Forest,
        Gradient Boosted Trees) and log each as a separate MLflow run
        with hyperparameters, AUROC, F1, and training duration.
  \item Use the MLflow Tracking UI to compare the three runs. Identify
        the best model by AUROC.
  \item Register the best model in the MLflow Model Registry with the
        name \texttt{payment-fraud-v2}. Transition it to the
        \texttt{Staging} stage.
  \item Load the staged model using \texttt{mlflow.spark.load\_model()}
        and run predictions on a 10{,}000-row test set. Verify the
        AUROC matches the training-time AUROC.
\end{enumerate}

\pe \textbf{Apache Beam: Windowed Transaction Aggregation.}
Implement an Apache Beam pipeline that aggregates payment transaction events:
\begin{enumerate}[label=(\alph*)]
  \item Generate a synthetic stream of 1 million transaction events with
        \texttt{(txn\_id, region, amount\_usd, event\_time)}. Introduce
        random event-time lag of 0--30 seconds to simulate mobile network delay.
  \item Write the pipeline using the Direct Runner (local execution).
        Apply 5-minute tumbling event-time windows and sum
        \texttt{amount\_usd} per \texttt{region} per window.
  \item Configure a composite trigger: emit an early result every 10 seconds
        of processing time AND emit a final result when the watermark passes
        the window end. Use accumulating mode.
  \item Set \texttt{allowed\_lateness} to 60 seconds. Verify that events
        arriving within 60 seconds after the window end are included in
        the final result.
  \item Compare the output of the 5-minute tumbling window against a
        10-minute sliding window with a 5-minute slide. Which produces
        more output records and why?
\end{enumerate}

%%====================================================================
\section*{Further Reading}
\label{sec:ch11-further}
%%====================================================================

\begin{itemize}
  \item \textbf{\cite{DeltaLake2020}} --- Open access on the VLDB
        website. Read Sections 2 (motivation), 3 (transaction log design),
        and 5 (performance evaluation). Section 4 (MERGE INTO) is worth
        reading for the upsert use case.
  \item \textbf{\cite{Lakehouse2021}} --- Four pages; reads quickly.
        Section 3 (Lakehouse Architecture) and Section 4 (Performance
        Evaluation) are most relevant to the course content. The final
        paragraph on DuckDB is prescient: by 2024, DuckDB had become
        a primary tool for in-process analytical queries on Delta tables.
  \item \textbf{\cite{Dataflow2015}} --- The most technically dense
        paper in the course. Read Section 1 (introduction) and Section 3
        (The Dataflow Model) carefully; Sections 4--6 describe the
        Google-internal execution system (MillWheel) which is less relevant
        to open-source practitioners. The appendix examples (event-time
        windowing with watermarks) are highly recommended.
  \item \textbf{\cite{MLTechnicalDebt2015}} --- Five pages. Every
        ML practitioner should read this. Map each of the nine debt
        categories to a tool in the MLOps stack described in this chapter.
  \item \textbf{\cite{Dehghani2022}} --- The definitive Data Mesh book.
        Chapters 1--3 (the problem) and 7--9 (the data product) are
        most relevant; Chapters 10--14 cover implementation which is
        beyond this course's scope.
  \item \textbf{Designing Data-Intensive Applications (Kleppmann, 2017)},
        Chapter 11 (\textit{Stream Processing}). The clearest explanation
        of event-time vs.\ processing time and watermarks for a
        practitioner audience. Pairs perfectly with the Dataflow Model paper.
  \item \textbf{Delta Lake documentation at \texttt{delta.io}}: the
        Quickstart guide and the ``How Delta Lake Works'' deep-dive are
        the best practical complement to the VLDB paper. The
        Databricks Community Edition (free) allows running all code
        examples in this chapter without a local Spark installation.
\end{itemize}
